\section{Discussion}
\label{sec:discussion}

The results show that calendar construction, absolute model adequacy and multiplicity control affect different parts of the dependence analysis in different ways. The stress-inference conclusion is highly stable: although each stress definition contains several nominally significant cells, none survives Bonferroni or Benjamini--Hochberg correction in either calendar panel. The appropriate conclusion is therefore not that dependence was constant during COVID-19 or 2024, but that the observed finite-threshold stress--calm differences do not provide multiplicity-adjusted evidence strong enough to support individual changes under the prespecified design.

This result is narrower than some crisis-focused commodity studies, but it is not contradictory. \citet{qiao2023} report stronger lower- and upper-tail contagion across commodity futures after the onset of COVID-19, while other international studies document state-dependent commodity and currency relationships \citep{ignatieva2017,go2024,ning2025}. Those studies use different asset universes, model structures, estimands and inferential procedures. The present design asks whether ten Ghana-relevant bivariate pairs exhibit stress-period changes at three finite thresholds after adjustment for 60 simultaneous comparisons. Nominal signals can coexist with the absence of corrected discoveries.

The copula goodness-of-fit results reveal a different form of model risk. Gold--Brent, Gold--WTI and Brent--WTI reject all eight tested static families in both panels. An AIC-best family still exists for such pairs, but the absolute GOF evidence shows why a relative winner should not be confused with an adequate model. The finding is consistent with a broader literature in which commodity and currency dependence varies over time and across regimes \citep{ignatieva2017,ning2025}. It motivates richer dynamic, regime-switching, vine or nonparametric dependence models, but the present evidence does not identify which of those alternatives would be adequate.

The cocoa-related pairs are especially informative about observation design. Cocoa--Gold, Cocoa--Brent and Cocoa--WTI reject all eight tested families under business-day alignment but only a subset under common-observation alignment. This is consistent with the nonsynchronous-trading literature, which shows that the observation process can alter measured dependence even when the underlying economic relationship is unchanged \citep{scholes1977,lo1990,martens2001,schotman2006}. In the present data the two panels remain strongly related on shared dates, yet their higher-order dependence classifications differ. Calendar sensitivity is therefore not a small preprocessing detail; it is a source of specification uncertainty.

The commodity--cedi non-rejection pattern requires more caution than the commodity-only results. None of the eight tested families is rejected for any commodity--cedi pair in either panel, but the cedi marginal fails the adequacy gate in both constructions. Because copula inference depends on the marginal transformation, these non-rejections cannot establish that the tested copulas are correct. \citet{fantazzini2009} shows that marginal misspecification can materially affect copula parameters and risk measures. The present cedi evidence is therefore conditional on an unresolved marginal model.

The cedi limitation is substantive rather than merely numerical. Panel A forward filling increases the cedi's exact-zero share, but Panel B still contains 16.51\% exact zeros. The PIT-clipping audit shows that the identified three-way floor tie is too small to alter the reported tail memberships or GOF classifications. The larger problem is that none of the tested continuous marginal candidates adequately represents the zero-heavy daily process. The positive GPD shape estimates in the EVT diagnostics therefore describe the fitted reference residual process, not a model-free structural tail index for the cedi.

The findings also refine the Ghana-specific literature. Previous studies report asymmetric, nonlinear and horizon-dependent relationships between commodities and the exchange rate \citep{archer2022,barson2023,yeboah2025,atipaga2025,woode2025}. The present results do not overturn that evidence. Instead, they show that at the daily-source frequency, under the chosen finite-tail estimands and multiplicity controls, robust stress-period discoveries are absent even while static copula inadequacy remains pronounced for several commodity pairs. Economically meaningful relationships may therefore coexist with weak corrected evidence for particular daily tail contrasts.

Three implications follow for empirical international-finance work. First, crisis labels should not substitute for formal inference. Second, relative model selection and absolute model adequacy should be reported separately. Third, the construction of cross-market return calendars should be documented as part of the statistical model whenever assets trade on different schedules. These are methodological implications rather than policy prescriptions; the study does not estimate causal pass-through, fiscal effects, reserve consequences or hedging effectiveness.

Several limitations remain. The cedi marginal is unresolved; the dependence analysis is bivariate and static; neither calendar achieves perfect intraday synchronization; Panel B contains variable one- to six-day intervals; the stress results depend on the stated windows, thresholds, block length and multiplicity procedures; and the real-data GOF analysis does not include a separate power study. The analysis also does not make unverified claims about upstream continuous-futures roll rules or back adjustment. These limitations define the scope of the conclusions rather than invalidate the documented cross-panel comparisons.

Overall, the most defensible results are the ones that survive multiple layers of scrutiny. The absence of corrected stress discoveries survives both calendars and all three stress definitions. Complete eight-family rejection survives for Gold--Brent, Gold--WTI and Brent--WTI. The equivalent conclusion does not survive for the cocoa-related commodity pairs. Commodity--cedi non-rejection is stable but remains conditional on marginal inadequacy. This hierarchy of evidence is the central contribution of the study.
